% =============================================================================
% SECTION 4: THEORETICAL ANALYSIS
% =============================================================================

\section{Theoretical Analysis}
\label{sec:theoretical}

Building on the TypeOracle and its two gates described in the previous section, this
section explains three properties that govern when and why DCA-Trie works: structural
faithfulness is preserved by construction, the search space shrinks by a bounded
factor, and tighter constraints do not guarantee higher accuracy. The third property
is the paper's central finding.

% -----------------------------------------------------------------------------
\subsection{Structural Faithfulness}
\label{sec:theoretical:faithfulness}
% -----------------------------------------------------------------------------

Faithfulness is preserved by construction. The TypeOracle operates as a filter on the
set $\mathcal{P}(\mathcal{G}, \mathcal{E}_q, L)$ of all paths within $L$ hops of the
question entities $\mathcal{E}_q$: it inspects each path and either admits or rejects
it, but it never constructs new paths, introduces new triplets, or modifies existing
ones. The admitted set is therefore a subset of $\mathcal{P}$, and every triplet in
every admitted path is a member of $\mathcal{T}$. The range gate and type gate are
predicates over existing triplets; they cannot fabricate new ones.

This property holds regardless of which gates are used or how they are parameterised.
Any filter applied after BFS enumeration preserves faithfulness, because filtering
removes paths without introducing new tokens. In v1, the trie is built from the
admitted subset, so trie lookup returns only tokens that correspond to valid triplet
prefixes. In v2, the trie is expanded at entity commits by adding only triplets
$(e_t, r, e')$ that exist in $\mathcal{G}$ and pass both gates
(Equation~\ref{eq:v2_expansion}), so newly added tokens correspond to valid triplets
and the same lookup argument applies.

% -----------------------------------------------------------------------------
\subsection{Search-Space Reduction}
\label{sec:theoretical:reduction}
% -----------------------------------------------------------------------------

The search-space reduction is bounded by construction. Let $|\mathcal{P}|$ be the
number of BFS-enumerated paths and $|\mathcal{P}_{\text{admitted}}|$ the number
surviving the TypeOracle. The reduction ratio
$\rho = 1 - |\mathcal{P}_{\text{admitted}}| / |\mathcal{P}|$ is bounded below by GCR
(no filtering, $\rho = 0$) and above by one path per question entity, since at least
one gold path must survive for a question to be answerable. The bounds themselves are
trivial; the empirical reduction depends on the ontology's structure. On the controlled
WebQSP subset, the TypeOracle removes $\rho = 0.138$ of paths (13.8\% of the candidate
set). This is not a theoretical guarantee but an empirical observation grounded in the
Freebase ontology's type system.

At each decoding step, trie lookup costs $O(L)$ where $L$ is the path length. In v2,
gate evaluation at entity commits costs $O(|\mathcal{N}(e_t)|)$, with each gate
evaluation being $O(1)$ because Freebase type sets are small (typically 1 to 5 types
per entity). The total gate evaluations across all entity commits is at most
$|\mathcal{G}|$, the number of triplets in the graph, so the amortised cost across all
steps is $O(|\mathcal{G}|)$ in the worst case, the same order as BFS enumeration in v1
but distributed across decoding rather than concentrated at construction time.

% -----------------------------------------------------------------------------
\subsection{Why Tighter Constraints Do Not Improve Accuracy}
\label{sec:theoretical:nonmonotone}
% -----------------------------------------------------------------------------

The empirical result from Section~\ref{sec:intro:results} shows that reducing SIR from
1.0 to 0.138 does not improve Hits@1. This is the non-monotone result: constraint
tightness and answer accuracy are not monotonically related. The reason is that the
TypeOracle trades off precision against recall.

Let $R$ be the recall of the oracle (the fraction of gold paths admitted) and $P$ its
precision (the fraction of admitted paths that are relevant). Tightening the constraint
increases $P$ but may decrease $R$ if the oracle incorrectly removes the gold path.
Hits@1 depends on both: the LLM must find the gold path among the admitted paths
(requiring $R > 0$) and rank it first (requiring low noise from irrelevant paths,
which $P$ controls). A second mechanism operates through the beam search. With beam
width $k$, the beam retains the $k$ highest-scoring partial paths at each step. Fewer
admitted paths mean a smaller, more homogeneous beam; if the gold path receives a low
LLM probability at an early step, the redundant paths that GCR retains can keep it in
the beam long enough to recover, whereas DCA-Trie has removed that redundancy. The
combined effect is that Hits@1 is constrained by the weaker of $R$ and the beam's
ability to rank the gold path first, so increasing $P$ while decreasing $R$ can leave
Hits@1 unchanged or reduced.

Practically, this means optimising the TypeOracle's gates to maximise precision
(removing as many irrelevant paths as possible) does not guarantee better answers. The
oracle must balance precision against recall, and the optimal balance depends on the
LLM's ability to navigate the remaining search space. SIR measures precision; Hits@1
measures the outcome of the balance. The two metrics answer different questions, and
improving one does not imply improvement in the other.